\question[3]
Beschreibe, wie im Unterricht mit dem Datenbankmanagementsystem \texttt{SQLite} gearbeitet wurde.

Gehe dabei auf die folgenden Aspekte ein:

\begin{itemize}
\item Herstellen einer Datenbankverbindung
\item Erstellen von Tabellen
\item Einfügen von Daten
\item Abfragen von Daten
\item Löschen von Daten
\end{itemize}

\begin{solutionbox}{6cm}
Wir verwenden das Datenbankmanagementsystem (DBMS) SQLite. Der Zugriff auf die Datenbank erfolgt aus Python heraus über das Modul sqlite3, das eine Schnittstelle zu SQLite bereitstellt.

\begin{itemize}
\item Mit \texttt{sqlite3.connect()} wurde eine Verbindung zu einer Datenbank hergestellt.
Die Verbindung zur Datenbank gibt uns einen Cursor zurück, mit dem wir unsere SQL Anweisungen
 übermitteln.
\item Mit dem SQL-Befehl \texttt{CREATE TABLE} wurden Tabellen angelegt.
\item Mit \texttt{INSERT INTO} wurden Datensätze eingefügt.
\item Mit \texttt{SELECT} wurden Daten aus Tabellen abgefragt.
\item Mit \texttt{DELETE FROM} wurden Datensätze gelöscht.
\end{itemize}
\end{solutionbox}